\todo{Lecture 19}
\begin{proposition}[Conditions de Penrose]
\begin{enumerate}
\item
\begin{align*}
AA^{+}A & =PDQQ^{*}D^{+}P^{*}PDQ \\
 & =PDD^{+}DQ \\
 & =PDQ=A
\end{align*}
\item 
\begin{align*}
(AA^{+})^{*} & =(A^{+})^{*}A^{*} \\
 & =P(D^{+})^{*}QQ^{*}D^{*}P^{*} \\
 & =P(D^{+})^{*}D^{*}P^{*} \\
 & =P(DD^{+})^{*}P^{*} \\
 & =PDD^{+}P^{*} \\
 & =AA^{+}
\end{align*}
\begin{align*}
AA^{+} & =PDQQ^{*}D^{+}P^{*} \\
 & =PDD^{+}P^{*}  \\
 & =(AA^{+})^{*}
\end{align*}
\item
$$
A^{+}AA^{+}=A^{+}
$$
\item 
$$
(A^{+}A)^{^{*}}=A^{+}A
$$

\end{enumerate}
\end{proposition}
\begin{theorem}
Soit $A\in \mathbf{C}^{m\times n}$. Il existe exactement une matrice $X\in \mathbf{C}^{n\times m}$ telle que les conditions de Penrose soient satisfaites :
\begin{enumerate}
\item $AXA=A$
\item $(AX)^{*}=AX$
\item $XAX=X$
\item $(XA)^{*}=XA$.
\end{enumerate}
\end{theorem}
\begin{proof}
On a que $X=A^{+}$ satisfait ces quatre conditions. On montre l'unicite. Soient $X,Y\in \mathbf{C}^{n\times m}$ satisfaisant les conditions de Penrose. On vérifie que $X=Y$.
\begin{enumerate}
\item 
\begin{align*}
X & =XAX  \\
 & =XAYAX \\
 & =XAYAYAYAX \\
 & =(XA)^{*}(YA)^{*}Y(AY)^{*}(AX)^{*} \\
 & =A^{*}X^{*}A^{*}Y^{*}YY^{*}A^{*}X^{*}A^{*} \\
 & =(AXA)^{*}Y^{*}YY^{*}(AXA)^{*} \\
 & =A^{*}Y^{*}YY^{*}A^{*} \\
 & =(YA)^{*}Y(AY)^{*} \\
 & =YAYAY \\
 & =YAY \\
 & =Y
\end{align*}
\end{enumerate}
\end{proof}
\section{Systèmes d'equations}
\subsection{Problème de moindres carres}
Soit $A\in \mathbf{R}^{m\times n}$ et $b\in \mathbf{R}^{m}$. On cherche $\min_{x \in \mathbf{R}^{n}}\lVert Ax-b \rVert$. On veut trouver $x^{*}\in \mathbf{R}^{n}$ solution de cet problème tel que $\lVert x^{*} \rVert\le \lVert y \rVert$ pour tout $y\in \mathbf{R}^{n}$ qui est une solution. Si $y^{*}$ est une solution, l'ensemble solution $S=\{ y^{*}+z:z\in \ker(A) \}$, on cherche $p \in S$ tel que $\lVert p \rVert$ est minimale. L'ensemble $S\cap \{ x^{*}\in \mathbf{R}^{n}:\lVert x \rVert\le\lVert y^{*} \rVert \}$ est compact, donc il existe un minimum dans cet ensemble.

Par une decomposition en valeurs singuliers, on a $A=PDQ$. Soit $x \in \mathbf{R}^{n}$, sa norme $\lVert x \rVert^{2}$ est egal a 
$$
x^{\mathsf{T}}x=x^{\mathsf{T}}Q^{\mathsf{T}}Qx=(Qx)^{\mathsf{T}}Qx=\lVert Qx \rVert^{2} 
$$
\begin{remark}
Pour $Q\in \mathbf{R}^{n\times n}$ unitaire et tout $x \in \mathbf{R}^{n}$ : $\lVert x \rVert=\lVert Qx \rVert$.
\end{remark}
\begin{theorem}
Une solution optimale au problème de moindres carres est donne par
$$
x=A^{+}b.
$$
\end{theorem}
\begin{proof}
On considère $\lVert Ax-b \rVert$ :
\begin{align*}
\lVert Ax-b \rVert  & =\lVert PDQx-b \rVert  \\
 & =\lVert PDQx-PP^{*}b \rVert  \\
 & =\lVert DQx-P^{*}b \rVert, \quad y:= Qx \\
 & =\lVert Dy-P^{*}b \rVert \\
 & =\lVert \mathrm{diag}(\sigma_{1},\dots,\sigma _{r},0,\dots,0)y-P^{*}b \rVert 
\end{align*}
Les solutions sont $y_{1}=\frac{(P^{*}b)_{1}}{b_{1}}$, \dots, $y_{r}=\frac{(P^{*}b)_{r}}{b_{r}}$. Les composantes $y_{r+1},\dots,y_{n}$ son toutes nulles. On a une norme minimale pour
$$
y=D^{+}P^{*}b.
$$
On re-substitue $y$ pour obtenir $x=Q^{*}D^{+}P^{*}b=A^{+}b$.
\end{proof}
\begin{example}
Trouver la solution minimale du système 
$$
\underbrace{ \begin{pmatrix}
0 & -1.6 & 0.6 \\
0 & 1.2 & 0.8 \\
0 & 0 & 0 \\
0 & 0 & 0 
\end{pmatrix} }_{ A }x=\begin{pmatrix}
5 \\
7 \\
3 \\
-2
\end{pmatrix}
$$
On calcule la pseudo-inverse de $A$ :
$$
A^{+}=\begin{pmatrix}
0 & 0 & 0 & 0 \\
-0.4 & 0.3 & 0 & 0 \\
0.6 & 0.8 & 0 & 0
\end{pmatrix}
$$
donc
$$
A^{+}b=\begin{pmatrix}
0 \\
0.1 \\
8.6
\end{pmatrix}
$$
est une solution optimale a ce système.
\end{example}

\section{Le meilleur sous-espace approximatif}
Etant donnes $k\le n \in \mathbf{N}$ et $a_{1},\dots,a_{m}\in \mathbf{R}^{n}$, on veut trouver $H\le \mathbf{R}^{n}$ de $\dim(H)\le k$ minimisant $\sum _{i=1}^{n}d(a_{i},H)^{2}$ ou $d$ est une distance.

Supposons que $H=\mathrm{span}\{ u_{1},\dots,u_{k} \}$ ou les $u_{i}$ sont orthonormaux. Une projection $\pi$ de $a_{i}$ sur $H$ est définie par
$$
\pi(a_{i})=\sum_{j=1}^{k} \langle a_{i},u_{j} \rangle u_{j}
$$
donc, la distance est donnée par
$$
d(a_{i},H)^{2}=\lVert a_{i}-\pi(a_{i}) \rVert ^{2}.
$$
Par Pythagore on a
$$
\lVert a_{i} \rVert ^{2}=d(a_{i},H)^{2}+\lVert \pi(a_{i}) \rVert ^{2}
$$
donc
$$
d(a_{i},H)^{2}=\lVert a_{i} \rVert ^{2}-\lVert \pi(a_{i}) \rVert ^{2}.
$$
On veut trouver
$$
\max_{u_{1},\dots,u_{k}} \sum_{i=1}^{m} \left\lVert \sum_{j=1}^{k} (a_{i}^{\mathsf{T}}u_{j})u_{j} \right\rVert ^{2}.
$$